\section{Observabilidad del entrenamiento: el loss spike es solo un síntoma}

Entrenar no consiste en enviar un job y esperar varias semanas. Las hipótesis de investigación pueden quedar invalidadas por data corruption, optimizer instability, precision overflow, code regression, network degradation o checkpoint mismatch. La global loss por sí sola no basta para diagnosticar, por eso se necesitan cuatro telemetry planes sincronizados: model, optimizer, data y system.

\subsection{Cuatro tipos de telemetry y anomaly triage}

Un \term{loss spike} es un aumento repentino y anómalo de la loss de entrenamiento; puede recuperarse en poco tiempo o puede anunciar que la run ya no es válida. El model plane registra la loss y las gradient/activation statistics; el optimizer plane registra el learning rate y la moment/update norm; el data plane registra el origen de cada lote, la token mixture y la corrupción; el system plane registra utilization, network, storage y hardware errors.

\lecturefigure{06-training-observability.png}{El diagnóstico del entrenamiento requiere cuatro tipos de telemetry sincronizada: model, optimizer, data y system.}{Subtítulos locales 17:20--20:10; concepto redibujado.}

\begin{knowledgebox}{Lectura de la figura: la correlación va antes que la root cause}
Cuando aparece un loss spike, primero se alinea el momento de la anomalía con el lote de datos, la optimizer update, los collective retries y los device errors, y solo después se plantea una hipótesis. Si solo se mira la loss, el equipo podría reducir el learning rate por error; si solo se mira el hardware log, podría pasar por alto una inestabilidad numérica desencadenada por datos específicos. Toda la telemetry debe usar un mismo run/step ID y guardar la versión del código, los datos y la configuración.
\end{knowledgebox}

\begin{warningbox}{Un dashboard en verde no significa que el experimento sea válido}
GPU utilization, throughput y job status pueden verse todos normales mientras el objetivo de entrenamiento ya está distorsionado por una mask errónea, datos contaminados, label shift o un optimizer state dañado. Un frontier program necesita scientific invariants: held-out loss, gradient range, distribución de los datos, checkpoint checksum y periodic evaluation.
\end{warningbox}

\subsection{Checkpoint, replay y resume}

Un \term{checkpoint} guarda los parámetros, el optimizer state, el scheduler state, el random state y el data cursor necesarios para reanudar el entrenamiento; es distinto del activation checkpointing, que ahorra memoria de GPU recalculando dentro de un solo forward/backward. El \term{deterministic replay} reproduce la anomalía, en la medida de lo posible, con el mismo código, el mismo lote y el mismo estado, para determinar si se trata de una transient infrastructure fault o de un recipe defect.

\lecturefigure{07-checkpoint-recovery.png}{La recuperación del entrenamiento debe conservar su significado científico: detect, checkpoint, replay, diagnose y, después, resume o corregir.}{Subtítulos locales 17:20--19:10; concepto redibujado.}

\begin{knowledgebox}{Lectura de la figura: el punto de recuperación debe incluir el data cursor}
Guardar solo los model weights pierde los optimizer moments, el LR schedule, la aleatoriedad y la posición en los datos, y al reanudar la curva puede cambiar sin que se note. Si en el replay el spike vuelve a aparecer en el mismo lote, conviene revisar primero data/recipe; si desaparece al cambiar de hardware, puede tratarse de una transient fault; si no se puede reproducir, hay que conservar la incertidumbre y reforzar el monitoreo, en lugar de declarar resuelto el problema.
\end{knowledgebox}

\teachervoice{Mann describe la mentalidad on-call como «vigilar el entrenamiento como se cuida a un paciente»: una anomalía en la loss curve puede hacerse visible solo horas después, y hacer rollback implica perder un avance costoso. Lo importante en clase no es dramatizar las guardias, sino que el entrenamiento mismo se ha convertido en un sistema de producción que requiere SRE discipline.}

\subsection{Resumen de la sección}

La observabilidad del entrenamiento conecta los experimentos con modelos y la operación del sistema. El loss spike es el punto de entrada; solo los cuatro tipos de telemetry, un checkpoint completo y un controlled replay permiten convertir una anomaly en conocimiento que sirve para corregirla.
